\properlane{mercy}{Mercy practiced before the meeting}
\subsection{An admired virtue and a neglected person}
Chrysostom begins with the Gospel's choice of virgins. Virginity is a demanding and honored discipline. A person who has undertaken it might suppose that the hardest work has been done. The parable refuses that conclusion: people marked by this discipline can still be excluded from the wedding. Chrysostom identifies the lamps with virginity and the oil with humanity, almsgiving and help for those in need. His inference depends on the distance between the accomplishment that others admire and the neighbor whose need may remain unanswered.\footnote{John Chrysostom, \textit{Homilies on Matthew}, 78.1.}

The difficulty is not resolved by lowering the value of virginity. Chrysostom continues to honor it. His argument is that success in one demanding good does not absolve a person of another obligation. He compares the power of bodily appetite with attachment to money and asks why someone who has resisted the first should surrender to the second. The force of his rebuke lies in the disparity: a person has accepted great labor and then allowed a more ordinary attachment to spoil its fruit.

That reasoning reaches beyond the specific state of consecrated virginity through the question it asks about conduct. The assembly includes people whose fidelity requires sacrifice. An inference from Chrysostom's argument is that an established record of sacrifice cannot answer the separate question of whether they help a person now in need. One kind of good does not perform another kind of good on its behalf. The hungry person still needs food; an injured relationship may still require patient attention. The argument remains concrete because Chrysostom names assistance, resources and persons rather than treating readiness as an admired spiritual disposition.

Matthew's account says that the foolish took no oil with them, but later has them report lamps going out. Chrysostom describes them as having insufficient oil. This is his explanation of the narrative difficulty, not a different text to substitute for the Gospel. He understands the failure as inadequate mercy: some preparation or some apparent readiness can coexist with a refusal to use one's resources generously. The image's severity concerns the failure that is revealed, not a numerical calculation of how many gifts purchase admission.

The entrance prayer begins from the other side of need. The person who will be asked to help first approaches God as someone asking to be heard. Augustine's account of the suffering Christ makes it possible to pray that lament without pretending that Christian life has already escaped affliction. Read with Chrysostom's exhortation, this petition invites an assembly that asks God for compassionate hearing to attend patiently to a neighbor's plea, even when it is inconvenient.

\subsection{Why the generous do not share}
The refusal to share oil is the strongest objection to identifying oil with mercy. If the wise are distinguished by generosity, their answer seems to exhibit the opposite. The request is limited: the foolish ask for some oil, not all. Their need is urgent, and the wise possess what might relieve it. Chrysostom himself sets out these features. He does not make the request unreasonable in order to make the refusal easier to accept.

His answer is that the final encounter reveals a limit of what one person can do for another. The wise cannot transfer their own life of mercy to the foolish. Their refusal signifies inability at that moment, not a lack of willingness to help. Chrysostom compares this separation with the gulf in the account of the rich man and Lazarus. His instruction concerns the present life, when resources can be given and another person's need can still be relieved.\footnote{Chrysostom, \textit{Homilies on Matthew}, 78.1--2.}

The sellers then receive a surprising identification. For Chrysostom they are the poor. They are available here, in the life in which mercy is practiced, and should have been approached earlier. The language of buying belongs to his explanation of the parable. It describes the lasting consequence of giving, not a price charged by the poor for divine grace. To make it an exchange that obliges God would undo the account of gift given by Wisdom and the Mass's prayers.

The resulting urgency has a precise object. Chrysostom challenges those who plan to leave money to the poor only at death. They assume that a final interval will remain in which to issue instructions and put affairs in order. A sudden death can remove that interval. The point is broader than the prudence of making a testament: postponing mercy because one expects a manageable last opportunity is itself a failure to respond while opportunity exists.

This explanation does not authorize refusal of assistance to a living person. Its premise is the difference between present action and the final disclosure of a completed life. A sick or dying neighbor belongs among those who can still be accompanied, comforted and helped. Treating such a person as already beyond mercy would misapply Chrysostom's argument at exactly the place where his exhortation demands action. Nor does the passage authorize pronouncing the fate of a particular dead person. The warning examines the hearer's present postponement.

The Collect's freedom for service can be understood here as release from habits that make help perpetually inconvenient. Some obstacles are burdens for which one rightly seeks relief; others include the attachment to possessions that Chrysostom criticizes. The prayer concerns service to God in the life of the whole person. In this interpretation, asking for freedom includes asking to become available to the neighbor, with both practical means and a willingness to use them.

\subsection{Seeking begins with an approach from Wisdom}
The moral demand might seem to leave every person alone to create the necessary goodness. Wisdom 6 describes a different beginning. The seeker is met by Wisdom's prior approach. Bonaventure explains that her own light enables the search and that love and action belong to it. Read alongside Chrysostom, this means that the capacity for mercy is received and exercised. The seeker responds to a gift by doing something; prior generosity does not make the response unnecessary.\footnote{Bonaventure, \textit{Commentarius in librum Sapientiae}, 6:13--17 Vulgate, Quaracchi VI, pp. 147--149.}

The image of Wisdom at the door also gives the search a proximity that differs from a journey toward an inaccessible possession. The one who wakes to seek finds that she is already near. That nearness can encourage a person whose past failures make a fresh beginning seem impossible. The Gospel's final closure gives urgency to responding, while Wisdom's present approach gives a reason to respond with confidence.

The psalm's thirsty body and lifted hands make this response bodily. Augustine connects the lifted hands of prayer with works of mercy. The same hands that ask may distribute food, receive a stranger or attend a sufferer. He does not dissolve prayer into useful action; he explains how the hands belong to both. His account of bodily thirst reaching toward resurrection also prevents the needy body from becoming a disposable occasion for someone else's spiritual improvement. The body cared for now belongs to the promised future.\footnote{Augustine, \textit{Expositions on the Psalms}, Psalm 62.3--5, 10.}

The acclamation's unknown hour closes off the strategy of delaying this action until a forecast makes it urgent. In his exposition of Matthew 24, Chrysostom makes continuing service the consequence of ignorance. The trusted servant has a charge to fulfill while the master is absent. For the hearer, readiness is therefore tested in the use of what is already entrusted, including money, responsibility and opportunities to assist.\footnote{Chrysostom, \textit{Homilies on Matthew}, 77.1--3.}

\subsection{Welcoming the stranger and recognizing Christ}
Gregory's Emmaus homily develops a different step in the argument. The disciples talk about Jesus, and Jesus walks with them, but they fail to recognize him. Gregory describes both love and doubt within them. Christ explains the Scriptures, yet the recognition comes after they urge him to stay and receive him at their table. Their uncertainty about his identity does not prevent an act of hospitality.

Gregory asks whether those who cannot yet love the traveler as God will at least love him as a stranger. This question gives the act its theological force. They do not first obtain certainty and then offer hospitality as an appropriate response to a distinguished guest. They welcome someone still unrecognized. Gregory understands their doing as a way into understanding the word they have heard. He then connects receiving the stranger with Christ's identification of himself among those whom the faithful receive.\footnote{Gregory the Great, \textit{Homiliae in Evangelia}, 23.1--2. The comparison with Matthew 25:35, 40 occurs in §2.}

This bears on Chrysostom's account of preparation without requiring Gregory to give the oil the same interpretation. The shared consequence is that the relation to Christ is exercised through a person whose needs can be overlooked. Chrysostom asks whether there has been mercy before the final arrival; Gregory explains how hospitality accompanies recognition of Christ now. The Luke Communion option recalls that recognition. It gives a particular reason for connecting reception at the Eucharistic table with receiving a person outside the celebration.

Where the psalm Communion option is used, the emphasis begins instead with Christ's provision. Augustine's shepherd feeds and refreshes the Church. The inference for mercy is that one who receives care from Christ can become willing to provide care to others. Bellarmine's qualification concerning hardship matters here: gratitude for the shepherd's care does not require someone to claim that all material need has disappeared. A poor person can receive and offer care without being recast as materially prosperous.\footnote{Augustine, \textit{Expositions on the Psalms}, Psalm 22.1--3; Bellarmine, \textit{Commentary on the Psalms}, O'Sullivan abridgment, Psalm 22.}

The Prayer over the Offerings concerns the Passion celebrated and inwardly received. In this interpretation, receiving Christ's self-giving calls for help that accepts inconvenience or cost for another person's good, beyond an easy act performed for approval.

\subsection{A meeting from which death does not exclude}
The common sleep of the virgins prevents this account of preparation from becoming a demand for endless exertion. Chrysostom reads the delay as contrary to an expectation of immediate arrival and the sleep as death. He connects the midnight cry with Paul's description of the Lord's descent. He also permits a literal nocturnal interpretation of the resurrection, whereas Augustine emphasizes unexpectedness. Neither interpretation supplies a calculable future date.\footnote{Chrysostom, \textit{Homilies on Matthew}, 78.1; Augustine, Sermon 93, §§6--9.}

The long second reading develops the consolation implicit in that comparison. Death does not deprive the faithful dead of their meeting with Christ. Aquinas explains why the living do not take precedence in a way that excludes them: the Lord's resurrection grounds the common hope. The gathering is therefore a reason to comfort a mourner, not a reason to demand that the mourner become emotionally invulnerable. Offering such consolation is itself a form of care, and it requires attention to the person's loss.\footnote{Aquinas, \textit{Super I ad Thessalonicenses}, 4, lectio 2.}

The short form grounds hope in Jesus' death and rising. The fuller scene of the gathering, including the trumpet and the order of rising, belongs to Paul's continuation. The comfort is grounded in Christ before it is developed through the scene of his coming.

The final prayer asks the Spirit to preserve the recipient's integrity. The need for help continues after the sacrament has been received and after an act of mercy has been performed. Perseverance includes renewing the response when another person's need is less attractive or when earlier generosity has brought no gratitude. The wedding remains the promised encounter with Christ; the continuing petition acknowledges dependence on the Spirit along the way.

\begin{fourSenses}
\item[Literal.] Wisdom encourages a receptive search, the psalm voices embodied desire, Paul gives hope concerning the dead, and Jesus warns that preparation cannot be improvised at the wedding's arrival. The prayers ask for hearing, freedom for service and preservation after reception.
\item[Allegorical.] Christ is the Bridegroom, risen Lord and shepherd. Gregory's Emmaus exposition identifies the Lord encountered through hospitality and recognized in the breaking of bread. Those who await Christ's coming belong to the Church he gathers.
\item[Moral.] Chrysostom's identification of the oil with mercy directs attention to help withheld or postponed. Gregory explains how acting on the word through hospitality can lead to understanding. Consoling a mourner belongs to the same practical care, with room for real sorrow.
\item[Anagogical.] The final meeting brings judgment and the hope of enduring life with the Lord. Death cannot exclude the faithful dead from his gathering; outward religious distinction cannot replace a life of mercy. The Spirit's help is sought for perseverance toward that meeting.
\end{fourSenses}
